% ============================================================
%  第12章 WebCL：使用OpenCL加速Web应用
%  源文件：content/chapter12/
% ============================================================


% ---- 12.1-chinese.md ----
\section{介绍WebCL}

Web应用如今普遍存在于Web浏览器中，并且云服务已成为我们访问个人和专业数据的主要方式。网页应用无需进行散布和安装就能在数亿个客户端电脑上运行。因为网页应用的跨平台特性，对于开发者非常具有吸引力。这也需要感谢互联网的快速发展，随着接连不断的移动链接，终端用户可以通过应用在任何地点、任何设备上获取自己感兴趣的信息。

随着图形处理器的计算能力逐渐增强，特别是在移动端，新一代的应用编程接口(比如：WebGL和WebCL)也运应而生。这些接口可以使用设备本身所携带的硬件加速器，让图像渲染和计算在用户端更加高效，并且可以延长电池寿命。这样的API不仅在有网络连接时，不会依赖于网络信号的质量，也可以在断网的情况下对网络应用进行加速。感谢WebGL和WebCL，让Web浏览器的性能达到了一个新的高度，同时也为网络应用开辟了一个新的应用领域。



% ---- 12.2-chinese.md ----
\section{如何使用WebCL编程}

WebCL 1.0其实就是使用JavaScript实现的OpenCL 1.2。同样的OpenCL API、语法以及运行时，无需太多行代码就能完成(JavaScript是面向对象的语言)。我们通常会将WebCL与其他类似的技术联系在一起进行比较。如果你对WebGL很不了解，也没有关系。

与OpenCL一样，WebCL的编程也是由两部分构成：

\begin{enumerate}
\item 主机端(例如，Web浏览器)用来控制和执行JavaScript程序
\item 设备端(例如，GPU)用来进行计算——OpenCL内核
\end{enumerate}

和WebGL一样，WebCL也属于窗口对象的特性。首先，我们要检查WebCL是否可用，如果可用，则创建对应的上下文对象：

\begin{codebox}[javascript]
// First check if the WebCL extension is installed at all
if (window.webcl == undefined){
  alert("Unfortunately your system does not support WebCL." +
  "Make sure that you have both the OpenCL dirver" +
  "and the WebCL browser extension installed.");
}

// Get a list of available CL platforms, and another list of the
// available devices on each platform. If there are no platforms
// or no available devices on any platform, then we can conclude
// that WebCL is not available
webcl = window.webcl
try{
  var platforms = webcl.getPlatforms();
  var devices = [];
  for (var i in platforms){
    var p = platforms[i];
    devices[i] = p.getDevices();
  }

  alert("Excellent! Your system does support WebCL");
} catch(e){
  alert("Unfortunately platform or device inquiry failed.");
}

// Setup WebCL context using the default device
var ctx = webcl.createContext();
\end{codebox}

\myGraphic{0.92}{images/chapter12/12-1.png}{图12.1 WebCL对象}

图12.1中描述了WebCL中的对象结构关系。注意WebCL是接口的名字，而webcl是JavaScript的对象。更多信息可以查阅WebCL的标准[1]。

Web应用无法获取平台信息之后，列出所有可以用的设备；而是通过获取到的设备信息了解对应平台，从而指定对应的设备或平台。

\begin{codebox}[javascript]
// Find appropriate device
for (var j = 0, jl = device.length; j < jl; ++j){
  var d = devices[j];
  var devExts = d.getInfo(cl.DEVICE_EXTENSIONS);
  var devGMem = d.getInfo(cl.DEVICE_GLOBAL_MEM_SIZE);
  var devLMem = d.getInfo(cl.DEVICE_LOCAL_MEM_SIZE);
  var devCompUnits = d.getInfo(cl.DEVICE_MAX_COMPUTE_UNITS);
  var devHasImage = d.getInfo(cl.DEVICE_IMAGE_SUPPORT);

// select device that matches your requirements
  platform = ...
  device = ...
}

// assuming we found the best device, we can create the context
var context = webcl.createContext(platform, device);
\end{codebox}

应用会在运行时对OpenCL对象进行管理，比如：命令队列、内存对象、程序对象、内核对象，以及入队命令(比如：执行内核、读取内存对象或写入内存对象)。

WebCL中定义了如下对象：

\begin{itemize}
\item \textbf{命令队列}
\item 内存对象(\textbf{数组}和\textbf{图像})
\item \textbf{采样器}对象，其描述了在内核中如何对图像进行读取
\item \textbf{程序}对象，其包含了一些列内核函数
\item \textbf{内核}对象，在内核源码中使用\_\allowbreak{}\_\allowbreak{}kernel声明的函数，其作为真正的执行对象
\item \textbf{事件}对象，其用来追踪命令执行状态，以及对一个命令进行性能分析
\item 命令同步对象，比如\textbf{标记}和\textbf{栅栏}
\end{itemize}

首先我们需要创建程序对象。WebCL与WebGL 1.0类似，假设可以提供一段内核源码。这样的话，Web应用就需要内嵌一个编译器。源码先从设备上进行载入，之后进行编译。和其他编译器一样，OpenCL编译器也定义了一些编译选项。

\begin{codebox}[javascript]
// Create the compute program from the source strings
program = context.createProgram(source);

// Build the program executable with relaxed math flag
try{
  program.build(device, "-cl-fast-relaxed-math");
} catch(err) {
  throw 'Error building program:' + err + program.getBuildInfo(device, cl.PROGRAM_BUILD_LOG);
}
\end{codebox}

如此这般，我们的程序就可以编译了，并且一个程序对象内具有一个或者多个内核函数。这些内核函数作为我们程序的入口函数，就如同静态库中的接口一样。为了区别不同的内核函数，我们需要创建WebCLKernel对象：

\begin{codebox}[javascript]
// Create the compute kernels from within the program
var kernel = program.createKernel("kernel_function_name");
\end{codebox}

就像普通函数一样，内核函数通常都会有一些参数。JavaScript会提供一些数据类型，typed arrays[2]就是用来传递不同类型的内核参数(具体参数类型见表12.1)。对于其他类型的数据，我们可以使用WebCL对象：

\begin{itemize}
\item WebCLBuffer和WebCLImage，可以用来包装数组
\item WebCLSampler可以对图像进行采样
\end{itemize}

一个WebCLBuffer对象可以将数据以一维的方式存储起来。数组中错存储的元素类型都可以使标量类型(比如：int或float)，向量类型，或是自定义结构体类型。

表12.1 setArg()中使用的webcl.type与C类型之间的关系

\par\medskip\noindent
\begingroup\fontsize{9bp}{12.5bp}\selectfont\renewcommand{\arraystretch}{1.12}
\begin{center}
\begin{tabularx}{\textwidth}{llXl}
\toprule
内核参数类型 & setArg()值的类型 & setArg()数组类型 & 注意 \\
\midrule
char, uchar & scalar & Uint8Array, Int8Arrary & 1 byte \\
short, ushort & scalar & Uint16Array, Int16Array & 2 bytes \\
int, uint & scalar & Uint32Array, Int32Array & 4 bytes \\
long, ulong & scalar & Uint64Array, Int64Array & 8 bytes \\
float & scalar & Float32Array & 4 bytes \\
charN & vector & Int8Array for (u)charN & N = 2,3,4,8,16 \\
shortN & vector & Int16Array for (u)shortN & N = 2,3,4,8,16 \\
intN & vector & Int32Array for (u)intN & N = 2,3,4,8,16 \\
floatN & vector & Float32Array for floatN and halfN & N = 2,3,4,8,16 \\
doubleN & vector & Float64Array for (u)doubleN & N = 2,3,4,8,16 \\
char, ..., double * & WebCLBuffer &  &  \\
image2d\_t & WebCLImage &  &  \\
sampler\_t & WebCLSampler &  &  \\
\_\_local &  & Int32Array([size\_in\_bytes]) & 内核内部定义大小 \\
\bottomrule
\end{tabularx}
\end{center}
\endgroup
\medskip

\begin{codebox}[javascript]
// Create a 1D buffer
var buffer = webcl.createBuffer(flags, sizeInBytes, optional srcBuffer);

// flags:
// webcl.MEM_READ_WRITE Default. Memory object is read and written by kernel
// webcl.MEM_WRITE_ONLY Memory object only writeten by kernel
// webcl.MEM_READ_ONLY Memory object only read by kernel
// webcl.MEM_USE_HOST_PTR Implementation requests OpenCL to allocate host memory
// webcl.MEM_COPY_HOST_PTR Implementation requests OpenCL to allocate host memory and copy data from srcBuffer memory. srcBuffer must be specified
\end{codebox}

注意，只有从内存对象或其子对象中读取数据，或读取多个具有重叠区域的子内存对象是有定义的。其他方式的并发读写都会产生未定义的行为。

WebCL图像对象可以存储1,2,3维的纹理，渲染内存或图像。图像对象中的元素可以从预定义的图像格式列表中选择。当前的WebCL版本最多只支持到二维图像。

\begin{codebox}[javascript]
// create a 32-bit RGBA WebCLImage object
// first, we define the format of the image
var imageFormat = {
  // memory layout in which pixel data channels are stored in the image
  'channelOrder':webcl.RGBA,
  // type of the channel data
  'channelType':webcl.UNSIGNED_INT8,
  // image size
  'width':image_width,
  'height':image_height,
  // scan-line pitch in bytes.
  // If imageBuffer is null, which is the default if rowPitch is not specified.
  'rowPitch'：image_pitch
};

// Image on device
// imageBuffer is a typed array that contain the image data already allocated by the application
// imageBuffer.byteLength >= rowPitch * image_height. The size of each element in bytes must be a power of 2.

var image = context.createImage(webcl.MEM_READ_ONLY | webcl.MEM_USE_HOST_PTR, imageFormat, imageBuffer);
\end{codebox}

WebCLSampler告诉内核函数读和对图像数据进行读取。WebCL的采样器和WebGL的采样器类似。

\begin{codebox}[javascript]
// create a smpler object
var sampler = context.createSampler(normalizedCoords, addressingMode, filterMode);
// normalizedCoods indicates if image coordinates specified are normalized.
// addressingMode indicated how out-of-range image coordinations are handled when reading an image.
// This can be set to webcl.ADDRESS_MIRRORED_REPEAT
// webcl.ADDRESS_REPEAT, webcl.ADDRESS_CLAMP_TO_EDGE,
// webcl.ADDRESS_CLAMP and webcl.ADDRESS_NONE.
// filterMode specifies the type of filter to apply when reading an image. This can be webcl.FILTER_NEAREST or webcl.FILTER_LINEAR
\end{codebox}

使用WebCLKernel.setArg()将标量、向量或内存对象以参数的形式传入内核。需要传递局部内存时，我们可以使用长度为1的Int32Array，其中的数值是需要分配多少字节的局部内存，因为局部内存不能在主机端或设备端初始化，不过主机端可以通过内核参数的形式告诉设备端要分配多少局部内存。

根据经验内核所需要的所有参数，都需要是一个对象。即使是标量也也要包装在长度为1的数组中。向量根据长度放置在数组中。数组、图像和采样器都是WebCL对象。内存对象(数组和图像)需要在内核执行命令入队前，通过主机内存转移到设备内存。

这里展示一个例子:

\begin{codebox}[javascript]
// Sets value of kernel argument idx with value as memory object or sampler
kernel.setArg(idx.a_buffer);
kernel.setArg(idx.a_image);
kernel.setArg(idx.a_sampler);

// Sets value of argument 0 to the integer value 5
kernel.setArg(0, new Int32Array([5]));

// Sets value of argument 1 to the float value 1.34
kernel.setArg(1, new Float32Array([1.34]));

// Sets value of argument 2 as a 3-float vector
// buffer should be  a Float32Array with 3 floats
kernel.setArg(2, new Float32Array([1.0, 2.0, 3.0]));

// Allocate 4096 bytes of local memory for argument 4
kernel.setArg(3, new Int32Array([4096]));
\end{codebox}

当需要传递一个标量时，需要告诉程序内核标量数据的类型。JavaScript只有一种类型——数字——所以我们需要通过setArg提供数据类型信息。

注意：

\begin{itemize}
\item 长整型是64位整数，其无法在JavaScript中表示。其只能表示成两个32位整数：低32位存储在数组的第一个元素中，高32位存储在第二个元素中。
\item 如果使用\_\allowbreak{}\_\allowbreak{}constant对内核参数进行修饰，那么其大小就不能超多webcl.DEVICE\_\allowbreak{}MAX\_\allowbreak{}CONSTANT\_\allowbreak{}BUFFER\_\allowbreak{}SIZE。
\item OpenCL允许通过数组传递自定义结构体，不过为了可移植性WebCL还不支持自定义结构体的传递。其中很重要的原因是因为主机端和设备端的存储模式可能不同(大端或小端)，其还需要开发者对于不同端的设备进行数据整理，即使主机和设备位于同一设备上。
\item 所有WebCL API都是线程安全的，除了kernel.setArg()。不过，kernel.setArg()在被不同的内核对象并发访问的时候也是安全的。未定义的行为会发生在多个线程调用同一个WebCLKernel对象时。
\end{itemize}

WebCL的对象中，命令队列中包含一系列操作和命令。应用可能使用多个独立的命令队列，直到其中有数据共享时才进行同步。

命令需要有序入队，不过在设备端可能是顺序执行(默认)，或者是乱序执行。乱序意味着当命令队列中包含两个命令A和B时，顺序的执行能保证A一定能在B之前执行，而乱序的话就不能保证A和B结束的顺序。对于乱序队列，可以使用等待事件或入队一个栅栏命令，来保证命令的之间的依赖顺序。乱序队列属于进阶的主题，本章就不再进行描述。对其感兴趣的读者可以去\href{http://sa09.ida.ucdavis.edu/docs/SA09-opencl-dg-events-stream.pdf}{这里}[3]看看。目前很多的OpenCL实现并不支持乱序队列。你可以尝试使用\texttt{QUEUE\_OUT\_OF\_ORDER\_EXEC\_MODE\_ENABLE}测试一下OpenCL实现是否支持乱序队列。如果得到的是\texttt{INVALID\_QUEUE\_PROERTIES}异常抛出的话，那么你所使用的设备就不支持乱序队列。

\begin{codebox}[javascript]
// Create an in-order command-queue(default)
var queue = context.createCommandQueue(device);

// Create an in-order command-queue with profiling of commands enabled
var queue = context.createCommandQueue(device, webcl.QUEUE_PROFILING_ENABLE)；

// Create an out-of-order command-queue
var queue = context.createCommandQueue(device, webcl.QUEUE_OUT_OF_ORDER_EXEC_MODE_ENABLE);
\end{codebox}

命令队列在指定设备上创建。多个命令队列具有其相对应的设备。应用可以通过主机端和设备端的数据传输，重叠的执行内核。图12.2展示了将问题分成两部分可能会更快的完成：

\myGraphic{0.92}{images/chapter12/12-2.png}{}

\begin{itemize}
\item 第一部分是主机端向设备端传输数据，分解之后只需要原来拷贝时间的一半即可完成工作。然后，执行内核，基本上也是一半的时间完成。最后将数据传回主机，同样还是原先一半的时间。
\item 第一个传输数据完成后，第二个才开始，有些类似于CPU的流水线。
\end{itemize}

当有一系列命令入队后，使用WebCL中的enqueueNDRange(kenrel, offsets, globals, locals)就能执行对应的命令:

\begin{itemize}
\item kernel——执行命令的内核对象。
\item offsets——对全局区域的偏移。如果传null，则代表offsets=[0, 0, 0]。
\item globals——内核所要解决问题的尺寸。
\item locals——每个维度上的工作组中，工作项的数量。如果传null，设备会对其自行设定。
\end{itemize}

例如，如果我们要操作一个宽度width，高度为height大小的图像，那么globals可以设置为[width, height]，locals是可以设置为[16, 16]。这里需要注意的是，如果enqueuNDRange()中locals的大小超过了webcl.KERNEL\_\allowbreak{}WORK\_\allowbreak{}GROUP\_\allowbreak{}SIZE，enqueuNDRange()就会执行失败。

\myRefRule

\begin{myRefs}

[1] WebCL 1.0 Specification, \texthref{http://www.khronos.org/registry/webcl/specs/latest/1.0}

[2] Typed Array Specification, \texthref{http://www.khronos.org/registry/typedarray/specs/latest}

[3] Derek Gerstmann Siggraph Asia 2009 on Advanced OpenCL Event Model Usage, \texthref{http://sa09.ida.ucdavis.edu/docs/SA09-opencl-dg-events-stream.pdf}

\end{myRefs}



% ---- 12.3-chinese.md ----
\section{同步机制}

与C/C++的OpenCL程序一样，设备可以以同步的方式对命令队列中的命令进行处理。主机端将命令提交到命令队列中，然后可以使用clFinish()等待命令队列上所有命令执行完成。

与之相似，在WebCL中webCLCommandQueue类具有两个参数：

\begin{enumerate}
\item event\_\allowbreak{}list——一个WebCLEevet数组
\item event——用来对设备执行命令的状态进行查询的事件对象
\end{enumerate}

通常，这两参数传的都是null。不过，当有事件传递给某个命令时，主机端可以使用clWaitForEvents用来等待某个命令执行结束。编程者也可以使用事件的回调函数，在对应命令完成时进行提示。这要求主机端代码在事件中提前注册一个回调函数。当event\_\allowbreak{}list传递如命令时，之后命令不会立即执行，只有等到event\_\allowbreak{}list上所有事件对象对应的命令执行完成，才会执行当前的命令。

\begin{codebox}[javascript]
// Enqueue kernel
try{
  kernel_event = new cl.WebCLEvent();
  queue.enqueueNDRange(kernel, 2, null, globals, locals, null, null, kernel_event);
} catch(ex) {
  throw "Couldn't enqueue the kenrel. " + ex;
}

// Set kernel event handling routines: call kernel_complete()
try{
  kernel_event.setCallback(webcl.COMPLETE, kernel_complete, "The kernel finished successfully");
} catch(ex){
  throw "Couldn't set callback for event. " + ex;
}

// Read the buffer
var data = new Float32Array(4096);
try{
  read_event = new webcl.WebCLEvent();
  queue.enqueueReadBuffer(clBuffer, false, 0, 4096 * 4, data, null, read_event);
} catch(ex){
  throw "Couldn't read the buffer. " + ex;
}

// register a callback on completion of read_event: calls read_complete()
read_event.setCallback(webcl.COMPLETE, read_complete, "Read complete");

// wait for both events to complete
queue.waitForEvents([kernel_event, read_event]);

// kernel callback
function kernel_complete(event, data){
  // event.status = webcl.COMPLETE or error if negative
  // event.data is null
  // data should contain "The kernel finished successfully"
}

// read buffer callback
function read_complete(event, data){
  // event.status = cl.COMPLETE or error if  negative
  // event.data contains a WebCLMemoryObject with values from device
  // data contains "Read complete"
}
\end{codebox}

以上代码所形成的应用，希望在命令结束时进行提示，那么首先需要创建一个WebCLEvent对象，将命令传递到事件中，然后注册一个JavaScript写的回调函数。注意WebCLEvent.setCallback()的最后一个参数可以是任意对象。还需要注意对WebCLBuffer、WebCLImage的数据读写，clBuffer的所有权从主机端传递到设备端。因此，当read\_\allowbreak{}complete()回调被调动时，clBuffer的所有权已经从设备端转移到主机端。所有权的转移意味着，当主机端没有对应内存的所有权时，是不能访问或使用对应的内存。当会回调函数被调用后，主机端就又能访问内存了。



% ---- 12.4-chinese.md ----
\section{WebCL的交互性}

再来说一下WebCL在计算方面的使用(而非渲染)。不过，数据已经在图形处理器处的时候，需要对数据进行渲染，是直接使用这段内存直接让OpenGL进行显式比较高效，还是将数据拷贝到CPU端，再让CPU交由OpenGL去显示高效呢？这里就需要WebGL的扩展能力了。

WebGL需要先创建上下文，然后再创建一个共享的WebCL上下文。这也就是WebGL能够将数据分享给WebCL的原因(流程见图12.3)，二者能够共享的内存对象如下：

\begin{itemize}
\item 纹理对象：包含图像的纹理信息
\item 顶点数组对象：包含顶点信息，比如坐标、颜色和法向量
\item 渲染数组对象：包含图像对象所用到的WebGL帧缓存对象[4]
\end{itemize}

\myGraphic{0.92}{images/chapter12/12-3.png}{图12.3 运行时调用WebCL和WebGL}

\myRefRule

\begin{myRefs}

[4] Iñigo Quilez. ShaderToy with Mandelbule shader, \texthref{http://www.iquilzles.org/apps/shadertoy/?p\$=\$mandelbulb}

\end{myRefs}



% ---- 12.5-chinese.md ----
\section{应用实例}

对于图像处理应用，通常会将轨迹跟踪输出到屏幕上。对于这样的应用，其表面都会使用多个三角形进行构成，然后通过WebGL进行渲染。计算内核代码提供了更加自由的方式，方便我们对通用计算进行优化。片元渲染器相较计算渲染器要难优化许多。更重要的是，纹理内存是缓存在处理器中，这样要比通过全局内存的方式访问CPU内存高效的多。不过，有些设备不支持图像内存，可以使用WebCLBuffer创建像素数组对象，从而更新WebGL的纹理。

本节中，我们使用Iñigo Quilez片元渲染器所渲染的影子玩具——曼德尔球[1]，将这个例子转化成一个WebCL内核，如图12.4所示。WebGL使用两个纹理三角形对画面进行填充。WebCL每帧都会生成新的纹理。因此，考虑到画布的大小(width，height)，WebCL将会生成width x height个像素点。

\myGraphic{0.92}{images/chapter12/12-4.png}{图12.4 WebGL中两个三角形可以画一幅WebCL生成的图像。}

当WebCL要使用WebGL内存进行计算时，WebGL上下文必须要进行初始化，并且在WebGL上下文上共享产生WebCL的上下文，然后就可以WebCL就能使用WebGL的相关内存对象进行计算。下面的代码中展示了，如何创建WebGL的纹理，并且将其用于WebCL进行计算：

\begin{codebox}[javascript]
// retrieva a <canvas> object with id glcanvas in HTML page
var canvas = documnet.getElementById("glcanvas");

// Try to grab the standard context. If it fails, fallback to experimental.
var gl = canvas.getContext("webcl") || canvas.getContext("experimental-webgl");

// Create OpenGL texture object
Texture = gl.createTexture();
gl.bindTexture(gl.TEXTURE_2D, Texture);
gl.texParameteri(gl.TEXTURE_2D, gl.TEXTURE_MAG_FILTER, gl.NEAREST);
gl.texParameteri(gl.TEXTURE_2D, gl.TEXTURE_MIN_FILTER, gl.NEAREST);

gl.texImage2D(gl.TEXTURE_2D, 0, gl.RGBA, TextureWidth, TextureHeight, 0, gl.RGBA, gl.UNSIGNED_BYTE, null);
gl.bindTexture(gl.TEXTURE_2D, null);

// Create OpenCL representation (a WebCLImage) of OpenGL texture
try{
  clTexture = context.createFormGLTexture2D(cl.MEM_WRITE_ONLY, gl.TEXTURE_2D, 0, Texture);
} catch(ex) {
  throw "Error: Failed to create WebCLImage." + ex;
}

// To use this texture, somewhere in your code, do as usual:
glBindTexture(gl.TEXTURE_2D, Texture);
\end{codebox}

将这个纹理作为参数传入内核。

\begin{codebox}[javascript]
kernel.setArg(0, clTexture);
kernel.setArg(1, new Uint32Array([TextureWidth]));
kernel.setArg(2, new Uint32Array([TextureHeight]));
\end{codebox}

最后，在内核代码中对WebCLImage对象进行调用。

\begin{codebox}[cpp]
__kernel
void compute(__write_only image2d_t pix, uint width, uint height){
  const int x = get_global_id(0);
  const int y = get_global_id(1);

  // compute pixel color as a float4
  write_imagef(pix, (in2)(x, y), color);
}
\end{codebox}

\myRefRule

\begin{myRefs}

[1] Iñigo Quilez。 ShaderToy with Mandelbulb shader, \texthref{http://www.iquilezles.org/apps/shadertoy/?p\$=\$mandelbulb}

\end{myRefs}



% ---- 12.6-chinese.md ----
\section{增强安全性}

安全性是各个Web浏览器公司非常重视的问题。你会将浏览器作为打开JavaScript脚本的默认程序吗？出于对未知程序的怀疑，其中可能包含有一些访问系统数据的操作，可能会从你的电脑中盗取数据，或是向你的电脑安装病毒或者木马。恶意脚本还会让计算机对资源的使用能力降低，甚至是让其崩溃。

WebGL和WebCL的例子中，其代码是为了使用GPU从而将性能最大化，不过当前的GPU缺少相应的硬件机制，一个程序中的相应内存不能被其他程序所访问到，这就需要内存保护和读取未初始化数据的机制。

保护机制的创建可以保证内核在WebCL中运行时，不会访问越界。WebCL中的内核可以以文本的形式传入，然后通过WebCL内置的编译器进行编译。当然，使用预编译文件的方式对于程序和数字签名处理来说更加安全，不过目前为止没有相应的内核二进制文件的格式标准。

两个能保护内存的准则：1. 分配的内存需要初始化(程序不能访问旧数据)；2. 访问合法内存的代码可以置为可信。

下面的机制可用来保证程序的安全性：

\begin{itemize}
\item 持续跟踪内存分配(如果平台支持动态分配，需要在运行时进行跟踪)
\item 监视合法读取的内存范围
\item 合法的数据砸编译和运行时，会显得更加高效
\end{itemize}

内核在传递时，需要使用一个校验器对其内容进行校验，以确保不存在访问越界的情况，不过这将影响(经测量，性能要比无校验的性能低30\%)。更多安全性的信息可以参阅[1]。

\myRefRule

\begin{myRefs}

[1] Vicit Ltd: Mikael Lepist, Tamiö Ylimäki, \texthref{http://learningwebcl.com/wp-content/uploads/2013/11/WebCLMemoryProtection.pdf}

\end{myRefs}



% ---- 12.7-chinese.md ----
\section{服务器端使用WebCL}

WebCL为计算而生，其可以使用在浏览器中，也可以使用与单独的JavaScript应用和服务中。Node.js是一款基于Chrome编译出的JavaScript运行时，其有运行速度快和网络应用可扩展能力强的特性。Node.js使用事件驱动，无阻塞输入/输出模型，让其集高性能和轻量化于一身。其实现的数据敏感、实时性的应用，在分布设备上的表现十分出色[1]。Node.js具有很好的模块化，其已经有超过80000个模块存在，其中一个就是\emph{node-webcl}[2]。\emph{node-webcl}也具有扩展特性，使用了一些Node.js标准中的特性(例如，使用Node.js内存时，和数组的方式一致)，其他特性在未来可能会考虑放在WebCL的标准中发布。

安装\emph{node-webcl}要比安装浏览器麻烦许多。其作为开发工具，需要依赖一些第三方库，并且需要重新编译。

\begin{itemize}
\item 首先，确认OpenCL的SDK能够使用你的CPU和GPU： - 对于Intel GPU，从\href{https://software.intel.com/en-us/vcsource/tools/opencl-sdk}{这里}下载SDK - 对于AMD GPU，从\href{https://developer.amd.com/tools-and-sdks/opencl-zone/amd-accelerated-paralled-processing-app-sdk}{这里}下载SDK - 对于NVIDIA GPU，从\href{https://developer.nvidia.com/opencl}{这里}下载SDK
\item 从\href{http://nodejs.org}{这里}安装Node.js
\item 安装node-gyp:npm install -g node-gyp
\item 安装node-image: - 在\href{http://freeimage.sourceforge.net}{这里}下载FreeImage库和头文件 - npm install node-image
\end{itemize}

如果你想使用WebGL的互操作扩展开发应用，你还需要安装node-webgl和node-glfw:

\begin{itemize}
\item 安装GLEW \href{http://glew.sourceforge.net}{地址}
\item 安装GLFW 3.x \href{http://www.glfw.org}{地址}
\item 安装AntTweakBar库和头文件 \href{http://anttweakbar.sourceforge.net}{地址}
\item npm install node-glfw node-webgl
\end{itemize}

建议将上面GLEW，AntTwekBar，FreeImage和GLFW包中的库文件和头文件，放置在你的头文件和库文件目录下面，这样在编译各个Node.js模块时，便于编译器进行查找；否则，你就需要修改binding.gyp文件。以上四个库在所有桌面平台(Mac，Windows和Linux)是可用的，并且在Linux和Mac下可以使用安装包管理器进行安装，比如apt-get和Homebrew[3]。

为了在代码中使用node-webcl，先打开编辑器，你需要将WebCL的API添加入全局命名空间中。下面webcl对象，就和浏览器中的window.webcl对象一样。

\begin{codebox}[javascript]
// add WebCL API
var webcl = require('node-webcl');

// rest of the code is identical to that in browser
\end{codebox}

将代码中的"require('node-webcl')"换成webcl=window.webcl就能让这段代码在支持WebCL的浏览器中生效。通过node-webcl，我们就能使用操作系统中的访问方式，并且能使用所有Node.js中的模块。

在Node.js中使用node-webcl，无论是对服务器上的计算进行加速，还是对浏览器的Web接口的交互进行加速[4]，对于网络访问都是十分有意义的事情[5]。

同样，不同的应用使用的编程框架也不同，比如:Python, R和MATLAB等。脚本语言提供一些特性，也可以用于JavaScript。虽然框架不同，但是很多运行时库也存在有JavaScript版本(甚至有工具能直接将其他脚本语言的代码，直接转换成JavaScript)。目前通过使用Node.js，JavaScript在现有的脚本语言框架中的性能是相对好的。

因为node-webcl的限制要比浏览器少，所以相关应用的开发进程通常会更快，并且性能还不错。由于Node.js的原因，也有很多动态和多线程应用也会使用JavaScript开发，不过这些应用可能不会在浏览器端使用。这些应用就是典型的服务器端应用，其需要在当前平台上的所有支持OpenCL的设备上，对复杂的负载进行调度。

随着Node.js和node-webcl的使用，基于JavaScript的应用可以快速的部署在服务器端和客户端。当前，一些数据敏感的实时应用，已经可以使用JavaScript完成。

目前，一个新的Node.js模块出现，其名为node-opencl[6]。其与node-webcl为同一个开发者，node-opencl是对node-webcl的重写。node-opencl要比node-webcl更加底层，其可以让JavaScript直接使用OpenCL的特性。WebCL只支持OpenCL 1.1，不过node-opencl支持所有版本的OpenCL。其他人可以使用JavaScript在node-opencl的基础上开发出自己的WebCL。不过，所有WebCL实现都会随着时间的推移，对特性进行增强，对Bug进行修复。所以，目前node-opencl和node-webcl都被Node.js所支持。

\myRefRule

\begin{myRefs}

[1] \texthref{https://www.nodejs.orh}

[2] \texthref{https://github.com/Motorola-Mobility/node-webcl}

[3] \texthref{http://brew.sh}

[4] \texthref{http://dev.w3.org/html5/websockets}

[5] \texthref{http://superconductor.github.io/superconductor}

[6] \texthref{https://github.com/mikeserven/node-opencl}

\end{myRefs}



% ---- 12.8-chinese.md ----
\section{WebCL的状态和特性}

WebCL 1.0已经在2014年3月14日发布，并且标准组织也已经准备好一致性测试用例[1]，白皮书和教程[2]。写下这段文字的时候，已经有4种WebCL 1.0的实现：

\begin{enumerate}
\item 诺基亚对火狐浏览器的WebCL实现[3]
\item 三星对基于WebKit浏览器的WebCL实现[4]
\item AMD对基于Chromium浏览器的WebCL实现[5]
\item 摩托罗拉移动对node.js(非浏览器)的WebCL实现[6][8]
\end{enumerate}

AMD开发的node-opencl服务器版本，使用node.js包装的OpenCL。[7]

一致性测试用例也在2014完成，这些测试例可以用来帮助Web浏览器供应商，检验其旗下浏览器对WebCL的支持程度。

\myRefRule

\begin{myRefs}

[1] \texthref{https://github.com/KhronosGroup/WebCL-conformance}

[2] \texthref{http://learningwebcl.com}

[3] \texthref{http://webcl.nokiaresearch.com}

[4] \texthref{https://github.com/SRA-SiliconValley/webkit-webcl}

[5] \texthref{https://github.com/amd/Chromium-WebCL}

[6] \texthref{https://github.com/Motorola-Mobility/node-webcl}

[7] \texthref{https://github.com/mikeseven/node-opencl}

[8] \texthref{https://github.com/mikeseven/node-webgl}

\end{myRefs}


